\documentclass[12pt]{article}
\usepackage{../../lecture_notes}
\usepackage{../../math}
\usepackage{../../uark_colors}

\title{Regression Discontinuity Design}
\author{Prof. Kyle Butts}
\date{}

\hypersetup{
  colorlinks = true,
  allcolors = ozark_mountains,
  breaklinks = true,
  bookmarksopen = true
}

\begin{document}
\maketitle
\begin{multicols}{2}

\section*{Motivation}

Many policies are typically decided according to some rules:
subsidies awarded based on income thesholds, minimum test-score cutoffs, DUIs given for being over legal-limit, or referendum pass by reaching at least $50\%$ vote share.

In these settings, being above or below the threshold will determine treatment status (or, at least partially).
In RDD parlance, we have a \emph{running variable} (or \emph{score}) $X_i$, a cutoff $c$, and a treatment indicator
$$
  D_i = \one{X_i \geq c}
$$
Units with $X_i < c$ are untreated, units with $X_i \geq c$ are treated.

Typically the running variable $X_i$ is correlated with a lot of unobservables, e.g. SAT test score is correlated with school smarts and family income.
Therefore comparing individuals above and below the threshold will be biased from selection into treatment. 

However, if we `zoom into' individuals quite near to the threshold, it's reasonable to think these individuals start to look very similar to each other. 
Is there really a difference between someone with a 1399 and a 1400 SAT score? 
As you move along the score variable, it is often reasonable to assume that characteristics evolve `smoothly'. 
If we think that, then outcomes should evolve smoothly too. 
Under this `continuity' of characteristics, the \emph{only thing} that is changing discretely right at the threshold is the treatment status.

If there is a treatment effect, then we should see a `jump' right at the threshold!
Under this `continuity' of characteristics idea, we can ascribe this jump as being due to treatment. 
This is the crux of the RDD identification strategy. 

Continuity, in general, feels very natural. 
Since most scores have some level of randomness to them (test scores, vote shares) or since thresholds are often arbitrary (round numbers for income threshhold or legal-limit), people tend to buy continuity.

However, there is a big problem you must think through. 
If people know the cut-off, then people could strategically try to \emph{manipulate} their score to end up on either side of the cut-off.
People might work slightly fewer hours to make sure they stay under an income threshold; certain kinds of students might retake the SAT test to reach a certain score; cops might `fudge the numbers' when recording the results of a breathalyzer test. 

These kinds of sorting behaviors create a problem for our `continuity' of characteristics assumption. 
Right at the cut-off, we might suspect a jump in characteristics due to sorting.
This would create a discontinuity at the cut-off in terms of $Y_i(0)$ that prevents us from calling the discontinuity in observed outcomes the causal effect of the policy.






% We distinguish between \emph{sharp} RDDs (treatment is a deterministic function of the score) and \emph{fuzzy} RDDs (treatment probability jumps but is not 0/1).

% \section*{Identification Intuitions}
% There are two complementary ways to think about identification in RDDs.
% \begin{enumerate}
%   \item \textbf{Local randomization.}  If the running variable contains enough idiosyncratic noise, then units with latent characteristics close to the threshold are effectively randomized into treatment or control by small perturbations.
%   In practice, this motivates comparing units in a narrow window $\left(c-h, c+h\right)$ and treating the design like a local experiment.
% 
%   \item \textbf{Continuity of potential outcomes.}  More commonly, we assume the conditional expectations of potential outcomes,
%   $$\mu_d(x) = \expec{Y_i(d)}{X_i = x},$$
%   are continuous in a neighborhood around $c$.
%   Under this continuity assumption, the only possible source of a jump in the observed outcome at $c$ is the treatment itself.
% \end{enumerate}

\section*{Identification Strategy}
The above discussion motivates out \emph{causal assumption} that the conditional expectations of potential outcomes,
$$\mu_d(x) = \expec{Y_i(d)}{X_i = x},$$
are continuous in a neighborhood around $c$.
Under this continuity assumption, the only possible source of a jump in the observed outcome at $c$ is the treatment itself.

The regression discontinuity estimand is the local average treatment effect at the cutoff:
\begin{align*}
  \tau_{\text{RD}} &= 
  \expec{Y_i(1) - Y_i(0)}{X_i = c} \\ 
  &= \lim_{x \downarrow c} \expec{Y_i}{X_i = x} - 
  \lim_{x \uparrow c} \expec{Y_i}{X_i = x}.
\end{align*}
In words: take the right-hand limit of the conditional mean and subtract the left-hand limit.

Under continuity, 
\begin{align*}
  \lim_{x \downarrow c} \expec{Y_i}{X_i = x} &=
  \lim_{x \downarrow c} \expec{Y_i(0)}{X_i = x} \\
  &= \expec{Y_i(0)}{X_i = c},
\end{align*}
where the last equality assumes continuity of potential outcomes (i.e. there is no `jump' at the cut-off).
Similarly, the right-hand limit equals $\expec{Y_i(1)}{X_i = c}$.

This identifies the conditional ATE for units with $X_i \approx c$.
When interpreting the estimate, remember that you only identify effects for people right at the cutoff (e.g. students with SAT score right around 1400).

\section*{When Continuity Fails}
Continuity can fail when agents can precisely manipulate the running variable or when other policies or treatments change at the same cutoff.
Examples include students retaking exams to clear a threshold, municipalities that ``game'' an index to qualify for a program, or spatial borders that feature many policies changing (e.g. town borders).
Such sorting contaminates one side of the cutoff and invalidates the simple continuity argument.

There are ways you can try and assess the plausability of the continuity assumption. 
First, you can ask ask whether units just above and below the cutoff are similar on observable pre-treatment characteristics — rerun the RD using prior outcomes or demographics as the outcome and test whether the estimated jump is zero.
Second, you can check the density of the running variable for mass points or spikes near the cutoff; a pile-up just above the threshold suggests strategic behavior (e.g., retaking an exam, rounding a score, or gaming an index).

\section*{Estimation Strategies}
There are several common estimators in practice.
The \emph{locally-constant} (difference-in-means) estimator compares average outcomes in $(c-h, c)$ and $(c, c+h)$.
The \emph{locally-linear} estimator fits simple linear models on each side and extrapolates to the cutoff; this is often implemented via the interacted regression
$$
Y_i = \alpha_0 + \alpha_1 D_i + \beta_0 (X_i - c) + \beta_1 D_i (X_i - c) + u_i
$$
on the subsample with $X_i \in (c-h, c+h)$, where recentering makes $\alpha_1$ the RDD estimate.
One can extend to higher-order local polynomials, but higher-order terms can induce extreme extrapolation and instability; prefer simpler local fits with careful bandwidth choice.

\section*{Bandwidth Selection}
Bandwidth selection balances a bias–variance trade-off: smaller $h$ reduces bias from extrapolation but increases variance; larger $h$ improves precision but risks misspecifying the local relationship.
Data-driven selectors (e.g. methods implemented in \texttt{rdrobust}) provide practical default choices, but applied work should report sensitivity across a range of bandwidths and polynomial orders.

\section*{Sorting and Falsification Checks}
Two standard gut checks are (i) covariate balance tests (apply the RD to pre-determined covariates) and (ii) density checks (McCrary test) for continuity of the running-variable distribution.
Additional falsifications include testing placebo cutoffs, testing leads/lags of the outcome, and checking whether the jump existed in pre-policy periods.

\section*{Extensions and Complications}

\subsection*{Spatial RDD}
Spatial RDDs view distance to a border as the running variable; they carry the same logic but raise additional concerns about geographic heterogeneity and multiple treatments at borders.
Multiple policies turning on at the same cutoff create bundled treatments and complicate interpretation.
Fuzzy RDDs occur when the probability of treatment jumps but is not deterministic; these are naturally interpreted with an instrumental-variables framework, estimating a local Wald ratio.

\subsection*{Difference-in-Discontinuities}
When a new policy is introduced into a context with pre-existing discontinuities, comparing RDs before and after the change (difference-in-discontinuities) can isolate the new policy effect under assumptions that pre-existing jumps are stable and that no new sorting occurs because of the new policy.

\end{multicols}

\section*{One-Page Checklist}
\begin{itemize}
  \item \textbf{Density:} Check for mass points or pile-ups in the running variable histogram; a spike just above the cutoff suggests strategic behavior or manipulation.
  
  Useful: \texttt{rddensity::rddensity()}.

  \item \textbf{Balance:} Test pre-treatment covariates and prior outcomes for jumps at the cutoff (e.g., prior grades, demographics, or baseline earnings).
  
  Useful: \texttt{rdrobust::rdplot()} and \texttt{rdrobust::rdrobust()}.

  \item \textbf{Context checks:} Ask whether other policies or treatments turn on at the cutoff.
  
  \item \textbf{Bandwidth:} Choose a baseline estimator (locally-linear is a good default), justify the window $h$, and report sensitivity to smaller and larger bandwidths.

  Useful: \texttt{rdrobust::rdrobust()} for automatic bandwidth selection and bias-corrected inference.
  
  \item \textbf{Functional form:} Favor simple local fits (linear) over high-order polynomials; where you compare polynomials, show how estimates change.

  Useful: local-polynomial options in \texttt{rdrobust} or by-hand fits with \texttt{fixest} on trimmed samples.

  \item \textbf{Plot:} Visualize the raw scatter, binned averages, and smoothed lines near the cutoff — look for a clear jump that is not driven by a few extreme bins.
  Useful: \texttt{rdrobust::rdplot()} or ggplot2 with binned summaries.
  
  \item \textbf{Report:} Present the main estimate with a sensitivity table/plot (varying bandwidths, kernels, and polynomial orders) and clearly document assumptions, diagnostics, and links to results from \texttt{rdrobust} and \texttt{rddensity}.
\end{itemize}

\end{document}
